\origpage{268--275}

\begin{center}
{\Large\sffamily\bfseries\color{BellaLavenderDeep} EXERCICES}
\end{center}
\vspace{0.8em}

\begin{enumerate}[label=\arabic*.,leftmargin=2.2em,itemsep=1.1em]

\item Soit $P\in\mathbb R[X]$. Existence et unicité de $Q\in\mathbb R[X]$ tel que $Q(0)=0$ et
\[
P(X)=Q(X)-Q(X-1).
\]

\item Calculer
\[
\prod_{k=1}^{n-1}\frac1{1-e^{2ik\pi/n}}.
\]

\item On pose
\[
\pi_n(z)=z^n+\frac1{z^n}\qquad\text{et}\qquad R(z)=z+\frac1z.
\]
Montrer qu'il existe un unique polynôme $P_n\in\mathbb Z[X]$ tel que, $\forall z\in\mathbb C^*$,
\[
\pi_n(z)=P_n(R(z)).
\]
Montrer que les racines de $P_n$ sont réelles, dans $[-2,2]$.

\item Trouver les polynômes $P$ de $\mathbb C[X]$ tels que
\[
P(X^2)=P(X)P(X+1).
\]

\item Quels sont les polynômes de $K[X]$ ayant une fonction polynôme associée constante, lorsque $K=\mathbb R$, puis $K=\mathbb Z/2\mathbb Z$.

\item Polynômes $P\in\mathbb R[X]$ tels que $X^6$ divise
\[
P(2X)-2(P(X))^2+1.
\]

\item Trouver les couples $(P,Q)\in(\mathbb R[X])^2$ tels que
\[
P^2=1+(X^2-1)Q^2.
\]

\item Condition sur $(p,q)\in\mathbb C^2$ pour que $X^8+X^4+1$ divise
\[
X^{8m}+pX^{4m}+q,
\]
$m$ fixé dans $\mathbb N^*$.

\item
\begin{enumerate}[label=\arabic*),leftmargin=2em]
\item Soit $P\in\mathbb Q[X]$ irréductible sur $\mathbb Q$. Montrer que ses racines (dans $\mathbb C$) sont simples.
\item Soit $P\in\mathbb Q[X]$, et $\lambda$ une racine dans $\mathbb C$, de multiplicité
\[
m(\lambda)>\frac12d^\circ P.
\]
Montrer que $\lambda\in\mathbb Q$.
\end{enumerate}

\item Soit $P\in\mathbb R[X]$ et $z_0\in\mathbb C$ avec $P(z_0)\ne0$ et $P(z)-P(z_0)$ qui admet $z_0$ comme racine d'ordre $\ge2$. Montrer qu'il existe $\eta>0$, tel que
\[
\forall\rho\in]0,\eta[,\quad
\operatorname{card}\{z\in\mathbb C,\ |z-z_0|=\rho,\ |P(z)|=|P(z_0)|\}\ge4.
\]

\item Soit $P\in\mathbb R[X]$, scindé de degré $n$,
\[
P(X)=\sum_{k=0}^n a_kX^k.
\]
Montrer que pour tout $k\in\{1,2,\ldots,n-1\}$,
\[
a_{k-1}a_{k+1}\le a_k^2.
\]

\item Soit une suite réelle $(a_n)_{n\in\mathbb N^*}$, avec $a_1>0$ et $a_p\ge0$ pour $p\ge2$. On lui associe les polynômes $P_n$, pour $n$ dans $\mathbb N^*$, définis par
\[
P_n(X)=\sum_{k=1}^n a_kX^k.
\]
\begin{enumerate}[label=\arabic*),leftmargin=2em]
\item Montrer que l'équation $P_n(X)=1$ admet une unique racine positive ou nulle notée $u_n$.
\item Étude de la suite $(u_n)$.
\item Dans le cas particulier où, $(\forall n\in\mathbb N^*)$, $(a_n=n)$, déterminer
\[
\lim_{n\to+\infty}u_n.
\]
\end{enumerate}

\end{enumerate}

\begin{center}
{\Large\sffamily\bfseries\color{BellaLavenderDeep} SOLUTIONS}
\end{center}
\vspace{0.8em}

\begin{enumerate}[label=\arabic*.,leftmargin=2.2em,itemsep=1.1em]

\item C'est un problème où l'inconnue, $Q$ intervient linéairement. On essaie de voir si $P$ est dans l'image d'un endomorphisme.

Soit $n=d^\circ P$, (si $P\ne0$). On vérifie que $Q$ doit être de degré $n+1$.

Soit
\[
E=\operatorname{Vect}(1,X,\ldots,X^n)\qquad\text{et}\qquad
F=\operatorname{Vect}(X,X^2,\ldots,X^{n+1}).
\]
Ces 2 espaces sont de dimension $n+1$, et $\theta:F\longrightarrow E$ définie par
\[
\theta(Q)(X)=Q(X)-Q(X-1)
\]
est injective, ($\theta(Q)=0$ et $Q\ne0\Rightarrow$, si $z$ est un zéro de $Q$, (dans $\mathbb C$), alors $\forall k\in\mathbb Z$, $z+k$ est zéro de $Q$ : absurde). Comme $\dim(E)=\dim(F)$, $\theta$ est bijective, donc pour $P\in E$, il existe une et une seule solution du problème dans $F$.

\item Les $e^{2ik\pi/n}$, $k=1,\ldots,n$, sont les zéros de $X^n-1$, sauf $x=1$, donc les zéros de
\[
\frac{X^n-1}{X-1}=1+X+X^2+\cdots+X^{n-1}=Q(X).
\]

Les $e^{2ik\pi/n}-1$ sont les zéros de $R(X)=Q(X+1)$, donc
\[
\prod_{k=1}^{n-1}(e^{2ik\pi/n}-1)
\]
est le produit des racines de
\[
R(X)=X^{n-1}+\cdots+n;
\]
ce produit vaut donc $(-1)^{n-1}n$. On a alors
\[
\prod_{k=1}^{n-1}\frac1{1-e^{2ik\pi/n}}=\frac1n.
\]

\item Le lien entre $z^n+\dfrac1{z^n}$ et $z+\dfrac1z$ vient de ce que
\[
\left(z+\frac1z\right)^n
=z^n+\frac1{z^n}+C_n^1\left(z^{n-2}+\frac1{z^{n-2}}\right)+C_n^2\left(z^{n-4}+\frac1{z^{n-4}}\right)+\cdots
\]
avec pour dernier terme
\[
\begin{cases}
C_{2p}^p,& n=2p,\\[2mm]
C_{2p+1}^p\left(z+\dfrac1z\right),& n=2p+1,
\end{cases}
\]
d'où l'on tire
\[
z^n+\frac1{z^n}
=(R(z))^n-C_n^1\pi_{n-2}(z)-C_n^2\pi_{n-4}(z)-\cdots-
\begin{cases}
C_{2p}^p,\\[1mm]
C_{2p+1}^p\pi_1(z).
\end{cases}
\]

Comme $\pi_1(z)=1\cdot R(z)=P_1(R(z))$ avec $P_1(X)=X\in\mathbb Z[X]$, et
\[
\pi_2(z)=z^2+\frac1{z^2}=\left(z+\frac1z\right)^2-2=P_2(R(z))
\]
avec $P_2(X)=X^2-2\in\mathbb Z[X]$ on justifie par récurrence, l'existence de $P_n(X)\in\mathbb Z[X]$, de degré $n$ tel que $\forall z\in\mathbb C^*$, $\pi_n(z)=P_n(R(z))$. De plus, $P_1$ et $P_2$ étant uniques, la relation de récurrence donne l'unicité de $P_n$, et le fait que $P_n$ soit de degré $n$.

Comme
\[
\pi_n(z)=\frac{z^{2n}+1}{z^n}
\]
s'annule en
\[
z_k=e^{i\pi/(2n)+ik\pi/n},\qquad k=0,1,\ldots,2n-1
\]
et que pour $k=0,1,\ldots,n-1$,
\[
R(z_k)=z_k+\overline z_k=2\cos\theta_k
\]
avec
\[
\theta_k=\frac\pi{2n}+\frac{k\pi}{n},
\]
$n$ valeurs distinctes de $[0,\pi]$, on a $n$ zéros distincts de $P_n$ de degré $n$ dans $[-2,2]$ : les racines de $P_n$ sont bien réelles et dans $[-2,2]$, ce sont les
\[
2\cos\left(\frac\pi{2n}+\frac{k\pi}{n}\right),\qquad k=0,1,\ldots,n-1.
\]

\item On est sur $\mathbb C[X]$ où les polynômes sont scindés, on peut utiliser les zéros.

Les polynômes constants solutions sont $P=0$ et $P=1$.

Si $P$ non constant est solution, soit $a$ un zéro de $P$, on a $P(a^2)=0$ donc $a^2$ zéro de $P$ et plus généralement, $\forall n\in\mathbb N^*$, $a^{2^n}$ est zéro de $P$. Donc si $a\ne0$ et $|a|\ne1$ on a déjà une infinité de $|a^{2^n}|$ distincts : absurde, d'où $a=0$ ou $|a|=1$.

Mais $x=a-1$ est aussi tel que $P(x^2)=0$, ... d'où $a-1=0$ ou $|a-1|=1$. Les 2 cercles
\[
\Gamma_1=\{z,|z|=1\}\qquad\text{et}\qquad\Gamma_2=\{z,|z-1|=1\}
\]
se coupent en
\[
a=\frac12+\frac{i\sqrt3}{2}\qquad\text{et}\qquad\frac12-\frac{i\sqrt3}{2}.
\]
Mais, avec $\varepsilon=1$ ou $-1$ et
\[
a=\frac12+\varepsilon\frac{i\sqrt3}{2}
\]
on a
\[
(a-1)^2=-\frac12-\varepsilon\frac{i\sqrt3}{2}\notin(\Gamma_1\cap\Gamma_2)\cup\{0,1\} :
\]
ces valeurs de $a$ sont exclues. Il reste $a\in\{0,1\}$.

Si on prend $P(X)=kX^p(X-1)^q$, la relation $P(X^2)=P(X)P(X+1)$ devient
\[
kX^{2p}(X^2-1)^q=k^2X^p(X-1)^q(X+1)^pX^q,
\]
vérifiée pour $k=1$ et $p=q$, ($k\ne0$ car on a supposé $P$ non nul).

Les solutions sont donc 1, 0 et les polynômes du type
\[
X^p(X-1)^p,\qquad p\in\mathbb N^*.
\]

\item Si $P$ de $\mathbb R[X]$, non constant, a une fonction polynôme $P$ de valeur constante $k$, le polynôme $P(X)-k$ admet plus de $n=d^\circ P$ zéros.

Soient $a_1,\ldots,a_{n+1}$, $n+1$ éléments distincts de $\mathbb R$, on aurait donc $P$ divisible par
\[
\prod_{j=1}^{n+1}(X-a_j),
\]
(Théorème 7.38), donc $d^\circ P>d^\circ P$ : absurde.

Seuls les polynômes constants de $\mathbb R[X]$ ont une fonction polynôme constante.

Si $K=\mathbb Z/2\mathbb Z$, les valeurs de la fonction polynôme associée à $P$ sont $P(0)$, $P(1)$ et le couple $(P(0),P(1))$ peut valoir $(0,0)$, $(0,1)$, $(1,0)$, $(1,1)$. Dans les 2 cas extrêmes, la fonction polynôme sera constante.

Or
\[
P(0)=P(1)=0\Longleftrightarrow P(X)=Q(X)X(X+1),
\]
(Théorème 7.38), et
\[
P(1)=P(0)=1\Longleftrightarrow P(X)=Q(X)X(X+1)+1,
\]
d'où les solutions du type
\[
P=X(X+1)Q+\lambda\qquad\text{avec }\lambda\in K=\mathbb Z/2\mathbb Z.
\]

\item En posant
\[
Q(X)=P(2X)-2(P(X))^2+1
\]
on a $X^6$ divise $Q$ si et seulement si
\[
Q(0)=Q'(0)=Q''(0)=\cdots=Q^{(5)}(0)=0.
\]

$Q(0)=P(0)-2(P(0))^2+1=0$ conduit à $P(0)=1$ ou $P(0)=-\dfrac12$. On trouve ensuite
\[
Q'(0)=P'(0)(2-4P(0))=0\Rightarrow P'(0)=0
\]
dans les 2 cas.

\[
Q''(X)=4P''(2X)-4P'^2(X)-4P(X)P''(X)
\]
conduit à
\[
P''(0)(1-P(0))=0.
\]

Si $P(0)=1\Rightarrow P''(0)=a\in\mathbb R$; $P(0)=-\dfrac12\Rightarrow P''(0)=0$.

\[
Q^{(3)}(X)=8P^{(3)}(2X)-12P'P''-4PP''',
\]
donc
\[
Q^{(3)}(0)=0=4P^{(3)}(0)(2-P(0))\Rightarrow P^{(3)}(0)=0.
\]

\[
Q^{(4)}(X)=16P^{(4)}(2X)-12(P'')^2-16P'P'''-4PP^{(4)},
\]
donc
\[
Q^{(4)}(0)=0=4P^{(4)}(0)(4-P(0))-12(P''(0))^2.
\]

Si $P(0)=1\Rightarrow P^{(4)}(0)=a^2$; si $P(0)=-\dfrac12\Rightarrow P^{(4)}(0)=0$.

Enfin
\[
Q^{(5)}(X)=32P^{(5)}(2X)-40P''P'''-20P'P^{(4)}-4PP^{(5)}
\]
donne
\[
Q^{(5)}(0)=4P^{(5)}(0)(8-P(0))=0\Longleftrightarrow P^{(5)}(0)=0.
\]

Les polynômes $P$ sont donc du type
\[
-\frac12+X^6R(X)
\]
ou
\[
1+a\frac{X^2}{2}+\frac{a^2X^4}{24}+X^6R(X)
\]
avec $R\in\mathbb R[X]$.

\item On a
\[
P\cdot P+(1-X^2)Q\cdot Q=1
\]
donc, (Bézout), $P$ et $Q$ sont premiers entre eux. Pour considérer les degrés on écarte le cas nul en l'examinant. Si $Q=0$, $P=1$ ou $P=-1$ donne des solutions. Mais $P=0$ ne donne pas de solution. Pour $Q\ne0$, on doit avoir $P\ne0$ et $d^\circ Q=d^\circ P-1$.

On dérive la relation :
\[
2PP'=2Q[XQ+(X^2-1)Q'],
\]
avec $P$ et $Q$ premiers entre eux donc $Q$ divise $P'$, et comme $d^\circ Q=d^\circ P'$, on a $\lambda$ réel tel que
\[
Q=\lambda P'.
\]
Alors
\[
P^2=1+\lambda^2(X^2-1)P'^2.
\]
Si $P$, (non nul) est de degré $n$ et de coefficient directeur $a$ en identifiant les termes de plus haut degré on obtient
\[
a^2=\lambda^2n^2a^2
\]
d'où
\[
\lambda=\pm\frac1n.
\]

Comme il est clair que si $(P,Q)$ est une solution, $(\varepsilon_1P,\varepsilon_2Q)$ avec $\varepsilon_1$ et $\varepsilon_2$ égaux à 1 ou $-1$ en est une autre, on suppose
\[
\lambda=\frac1n\qquad\text{et}\qquad Q=\frac{P'}n.
\]

On cherche $P$ polynôme de degré $n$, vérifiant
\[
P^2=1+(X^2-1)\frac{P'^2}{n^2},
\]
ou encore
\[
(1-x^2)P'^2-n^2(1-P^2)=0.
\]

L'équation différentielle
\[
\frac{y'}{\sqrt{1-y^2}}=\pm\frac n{\sqrt{1-x^2}}
\]
que doit vérifier la restriction, $y$, de $P$ à $]-1,1[$, conduit à
\[
\operatorname{Arcsin}y=\pm(n\operatorname{Arcsin}x+k)
\]
et à
\[
y=\pm\sin(n\operatorname{Arcsin}x+k)
\]
soit à
\[
y=\pm[\sin(n\operatorname{Arcsin}x)\cos k+\cos(n\operatorname{Arcsin}x)\sin k].
\]

Un calcul classique conduit à
\[
\sin n\theta=C_n^1\cos^{n-1}\theta\sin\theta-C_n^3\cos^{n-3}\theta\sin^3\theta+\cdots
\]
et à
\[
\cos n\theta=\cos^n\theta-C_n^2\cos^{n-2}\theta\sin^2\theta+\cdots
\]
Si $n$ est pair, les $\cos^{n-2k-1}(\operatorname{Arcsin}x)$ contiendront $\sqrt{1-x^2}$ en facteur, ce ne sera pas un polynôme, donc $\cos k$ doit être nul, de même $n$ impair impose $\sin k$ nul et on trouve :

si $n=2p$, $P(X)$ doit être du type $\pm\cos(2p\operatorname{Arcsin}x)$ sur $]-1,1[$

et si $n=2p+1$, $P(X)$ doit être du type $\pm\sin((2p+1)\operatorname{Arcsin}x)$ sur $]-1,1[$.

Il reste à vérifier que ces polynômes vérifient la relation du départ, (car on a dérivé), avec
\[
Q=\frac\varepsilon nP',\qquad(\varepsilon\in\{-1,1\}).
\]

Si par exemple
\[
P=\varepsilon_1\cos(2p\operatorname{Arcsin}x),
\]
alors
\[
Q=-\frac{\varepsilon\varepsilon_1}{2p}\frac{2p\sin(2p\operatorname{Arcsin}x)}{\sqrt{1-x^2}}
\]
et
\[
P^2+(1-x^2)Q^2=\cos^2(2p\operatorname{Arcsin}x)+\sin^2(2p\operatorname{Arcsin}x)=1.
\]
On vérifierait de même que la solution trouvée dans le cas $n=2p+1$ convient.

\item Une condition nécessaire et suffisante est que les racines de $X^8+X^4+1$ soient racines de $X^{8m}+pX^{4m}+q$ avec des multiplicités supérieures ou égales.

Or
\[
(1+X^4+X^8)(1-X^4)=1-X^{12} :
\]
les zéros cherchés sont simples, ce sont les racines $12^\text{ième}$ de l'unité, non racines $4^\text{ième}$, soit les
\[
e^{2ik\pi/12}
\]
avec $k=1,2;4,5;7,8;10,11$, on doit donc avoir
\[
e^{8m\,2ik\pi/12}+pe^{4m\,2ik\pi/12}+q=0
\]
pour ces valeurs de $k$. Avec
\[
j=e^{2i\pi/3},
\]
comme
\[
4m\cdot\frac{2ik\pi}{12}=mk\frac{2i\pi}{3}
\]
on doit avoir
\[
j^{2mk}+pj^{mk}+q=0
\]
pour $k=1,2,4,5,\ldots$

Il reste finalement les conditions
\[
j^{2m}+pj^m+q=0,\qquad\text{et}\qquad j^{4m}+pj^{2m}+q=0,
\]
car
\[
j^{8m}=j^{6m}j^{2m}=(j^3)^{2m}j^{2m}=j^{2m}
\]
et
\[
j^{4m}=j^{3m}j^m=j^m.
\]

\item
\begin{enumerate}[label=\arabic*),leftmargin=2em]
\item Le polynôme $P$, irréductible sur $\mathbb Q$ n'a pas de diviseurs dans $\mathbb Q[X]$, il est donc premier avec $P'$, (qu'il ne divise pas) donc $\exists U$ et $V$ dans $\mathbb Q[X]$ tels que
\[
UP+VP'=1.
\]

Comme $\mathbb Q[X]$ est plongé dans $\mathbb C[X]$, si $\lambda$ était racine, (dans $\mathbb C$), multiple de $P$, on aurait $P(\lambda)=P'(\lambda)=0$ d'où $1=0$ : exclu. Donc les zéros de $P$, dans $\mathbb C$, sont simples.

\item On suppose cette fois que $P$ admet, dans $\mathbb C$, un zéro $\lambda$ de multiplicité
\[
m(\lambda)>\frac12d^\circ P.
\]

L'idéal
\[
I=\{Q;Q\in\mathbb Q[X],Q(\lambda)=0\}
\]
est principal, et il contient $P$. Soit $A$ un générateur de $I$, $\lambda$ est racine simple de $A$, sinon $A'(\lambda)=0$, $A'\in\mathbb Q[X]$ et $d^\circ A'<d^\circ A$ avec $A'$ multiple de $A$, c'est absurde.

Avec $P_1\in\mathbb Q[X]$ tel que $P=AP_1$, on a $\lambda$ racine simple de $A$ et d'ordre $m(\lambda)$ de $P$ donc $\lambda$ racine d'ordre $m(\lambda)-1$ de $P_1$. Si $m(\lambda)>1$, alors $P_1\in I$ : il existe $P_2$ tel que $P_1=AP_2$ et $P=A^2P_2$. Si $m(\lambda)>2$, on a $\lambda$ zéro d'ordre $m(\lambda)-2$ de $P_2$ qui est donc multiple de $A$, on itère et finalement il existe un polynôme $P_{m(\lambda)}$ tel que
\[
P=A^{m(\lambda)}P_{m(\lambda)}.
\]

Si $\lambda\notin\mathbb Q$, $A$ ne peut pas être de degré 1, donc $d^\circ A\ge2\Rightarrow d^\circ(A^{m(\lambda)})\ge2m(\lambda)>d^\circ P$, c'est absurde.

Donc $\lambda\in\mathbb Q$.
\end{enumerate}

\item Si $P(z)-P(z_0)$ admet $z_0$ pour zéro d'ordre $q\ge2$, $P',P'',\ldots,P^{(q-1)}$ sont nuls en $z_0$, d'où par Taylor Young, l'existence d'un développement limité
\[
P(z_0+h)=P(z_0)+a_qh^q+o(h^q)
\]
avec $a_q\ne0$. On a alors
\[
\begin{aligned}
|P(z_0+h)|^2
&=(P(z_0)+a_qh^q+o(h^q))(\overline{P(z_0)}+\overline{a_q}\,\overline h^{\,q}+o(h^q))\\
&=|P(z_0)|^2+h^qa_q\overline{P(z_0)}+\overline h^{\,q}\overline{a_q}P(z_0)+o(h^q).
\end{aligned}
\]

Comme $a_q\overline{P(z_0)}\ne0$, posons
\[
a_q\overline{P(z_0)}=Ae^{i\alpha}
\]
avec $A>0$, et soit $h=\rho e^{i\theta}$, on a
\[
|P(z_0+\rho e^{i\theta})|^2=|P(z_0)|^2+\rho^qA\left(e^{i(q\theta+\alpha)}+e^{-i(q\theta+\alpha)}+o(1)\right).
\]

Avoir $|P(z_0+\rho e^{i\theta})|=|P(z_0)|$ équivaut, (pour $\rho\ne0$) à avoir $\theta$ tel que
\[
\rho^qA\left(e^{i(q\theta+\alpha)}+e^{-i(q\theta+\alpha)}+o(1)\right)=0.
\]

Soit
\[
\varphi(\theta)=2\cos(q\theta+\alpha)+o(1),
\]
soit $\varepsilon\in]0,1[$ fixé, il existe $\eta>0$ tel que
\[
\rho=|h|\in]0,\eta[\Rightarrow |o(1)|<\varepsilon<1.
\]
Mais alors comme
\[
\cos\left(q\cdot\frac{2k\pi-\alpha}{q}+\alpha\right)=\cos2k\pi=1
\]
et
\[
\cos\left(q\cdot\frac{2k\pi+\pi-\alpha}{q}+\alpha\right)=\cos(2k+1)\pi=-1,
\]
on a
\[
\varphi\left(\frac{2k\pi-\alpha}{q}\right)>0
\]
et
\[
\varphi\left(\frac{2k\pi+\pi-\alpha}{q}\right)<0,
\]
la fonction $\varphi$ s'annule sur les intervalles
\[
\left]\frac{n\pi-\alpha}{q},\frac{n\pi+\pi-\alpha}{q}\right[,
\]
ceci pour $0\le n<2q$ : on a $2q$ intervalles de ce type dans
\[
\left[-\frac\alpha q,2\pi-\frac\alpha q\right]
\]
d'où $2q\ge4$ éléments $z=z_0+h$ avec $|h|=\rho$ tels que $|P(z)|=|P(z_0)|$ et ce pour tout $\rho\in]0,\eta[$.

\item Comme
\[
P(X)=\sum_{k=0}^n\frac{P^{(k)}(0)}{k!}X^k,
\]
les $a_k$ sont liés aux $P^{(k)}(0)$. Soit $x_1,\ldots,x_r$ les zéros réels distincts de $P$, de multiplicités $\alpha_1,\ldots,\alpha_r$, indexés en croissant : $x_1<x_2<\cdots<x_r$. On a $\alpha_1+\alpha_2+\cdots+\alpha_r=n$, ($P$ scindé) puis $x_1,\ldots,x_r$ sont zéros d'ordre $\alpha_1-1,\ldots,\alpha_r-1$ de $P'$ et, par Rolle, entre $x_i$ et $x_{i+1}$, $P'$ s'annule, donc on a $y_1,\ldots,y_{r-1}$, soit $r-1$ autres zéros de $P'$ qui admet donc
\[
\sum_{i=1}^r(\alpha_i-1)+r-1=n-1
\]
zéros. On a $P'$ scindé, et par récurrence chaque polynôme $P^{(j)}$ est scindé.

Puis pour
\[
Q=\prod_{i=1}^p(x-z_i),
\]
scindé, (les $z_i$ distincts ou non), on a
\[
\frac{Q'}Q=\sum_{i=1}^p\frac1{x-z_i}
\]
donc, en dérivant,
\[
\frac{Q''Q-Q'^2}{Q^2}=-\sum_{i=1}^p\frac1{(x-z_i)^2}
\]
donc, pour $x\notin\{z_1,\ldots,z_p\}$ on a
\[
Q''Q-Q'^2\le0,
\]
résultat vrai aussi pour les $z_i$ par continuité.

On applique ceci pour le polynôme $Q=P^{(j-1)}$, cela donne
\[
P^{(j+1)}P^{(j-1)}-(P^{(j)})^2\le0,
\]
et en 0 on obtient
\[
(j+1)!a_{j+1}(j-1)!a_{j-1}\le(j!)^2(a_j)^2
\]
d'où
\[
a_{j-1}a_{j+1}\le\frac{(j!)^2}{(j-1)!(j+1)!}a_j^2
=\frac{j}{j+1}a_j^2\le a_j^2.
\]

\item On constate que
\[
P'_{n+1}(x)=a_1+\sum_{k=2}^{n+1}ka_kx^{k-1},
\]
donc $P'_{n+1}(x)\ge a_1>0$ sur $[0,+\infty[$ : la fonction $P_{n+1}$ est strictement croissante sur $[0,+\infty[$, avec $P_{n+1}(0)=0$ et
\[
\lim_{x\to+\infty}P_{n+1}(x)=+\infty,
\]
il existe donc un et un seul $u_{n+1}$ tel que $P_{n+1}(u_{n+1})=1$, et ce pour tout $n\ge1$. Comme $P_1(x)=a_1x$ est aussi croissant de 0 à $+\infty$, $u_1$ existe de même.

2) On a
\[
P_{n+1}(x)=P_n(x)+a_{n+1}x^{n+1}.
\]
Donc
\[
P_{n+1}(u_n)=P_n(u_n)+a_{n+1}u_n^{n+1}\ge P_n(u_n)=1,
\]
vu la croissance de $P_{n+1}$, c'est que $u_{n+1}\le u_n$.

La suite $(u_n)_{n\in\mathbb N}$, décroissante, minorée par 0 est donc convergente.

3) On a
\[
P_2(X)=X+2X^2,
\]
donc
\[
P_2(X)=1\Longleftrightarrow 2X^2+X-1=0
\Longleftrightarrow X=-1\text{ et }X=\frac12,
\]
on veut $u_2>0$, donc $u_2=\dfrac12$, et, par décroissance,
\[
\forall n\ge2,\qquad u_n\in\left[0,\frac12\right].
\]

Les polynômes $P_n$ avec
\[
P_n(X)=X+2X^2+3X^3+\cdots+nX^n
\]
font penser à la dérivée de la série entière de $\dfrac1{1-X}$.

Pour $|x|<1$ on a
\[
\sum_{n=0}^{\infty}x^n=\frac1{1-x}
\]
d'où par dérivation terme à terme des séries entières sur l'intervalle ouvert de convergence,
\[
\sum_{n=1}^{\infty}nx^{n-1}=\frac1{(1-x)^2}
\qquad\text{et}\qquad
\sum_{n=1}^{\infty}nx^n=\frac{x}{(1-x)^2}.
\]

Si on pose
\[
R_n(x)=\sum_{k=n+1}^{\infty}kx^k,
\]
on a donc la relation
\[
P_n(x)+R_n(x)=\frac{x}{(1-x)^2},
\]
et en particulier
\[
1+R_n(u_n)=\frac{u_n}{(1-u_n)^2}.
\]

La suite des fonctions $R_n$ converge uniformément vers 0, sur $\left[0,\dfrac12\right]$, les $u_n$, pour $n\ge2$ sont dans $\left[0,\dfrac12\right]$, donc
\[
\lim_{n\to+\infty}\frac{u_n}{(1-u_n)^2}=1
\]
et c'est aussi
\[
\frac{l}{(1-l)^2}
\]
en notant
\[
l=\lim_{n\to+\infty}u_n,
\]
qui existe, avec $l\in\left[0,\dfrac12\right]$.

L'équation
\[
l^2-3l+1=0
\]
donne
\[
l=\frac{3\pm\sqrt5}{2},
\]
mais la seule racine dans $\left[0,\dfrac12\right]$ est
\[
l=\frac{3-\sqrt5}{2},
\]
donc
\[
\lim_{n\to+\infty}u_n=\frac{3-\sqrt5}{2}.
\]

\end{enumerate}
